\chapter{Embeddings And Representations}

Tokenization gave TinyGPT a stream of integers. PyTorch can store those integers, and a data loader can batch them, but the integers themselves are not yet useful for language understanding. The number \(2\) is not automatically ``more like'' the number \(3\) than the number \(8\). Token IDs are labels, not measurements.

Embeddings are the first trainable bridge from symbolic text to numerical learning. They turn token IDs into vectors, and those vectors become the working memory of the model.

\section{Chapter Goals}

By the end of this chapter, you should be able to:

\begin{itemize}
  \item explain why token IDs are arbitrary labels;
  \item describe an embedding table as a trainable lookup matrix;
  \item trace TinyGPT token IDs into vectors with shape \shape{[B,T,C]};
  \item explain why GPT adds token embeddings and position embeddings;
  \item connect embedding updates to backpropagation and next-token loss;
  \item explain weight tying between input embeddings and the output head;
  \item calculate embedding parameter counts for tiny and larger models;
  \item read and debug the source code that creates initial GPT representations.
\end{itemize}

\section{The Running TinyGPT Example}

We continue with the TinyGPT sentence:

\begin{quote}
\texttt{Tiny GPT models learn from tokens .}
\end{quote}

One toy vocabulary is:

\begin{longtable}{p{0.16\textwidth}p{0.22\textwidth}p{0.46\textwidth}}
\toprule
\term{ID} & \term{Token} & \term{Comment} \\
\midrule
0 & \texttt{<pad>} & optional padding token \\
1 & \texttt{<unk>} & unknown token \\
2 & \texttt{Tiny} & first content token \\
3 & \texttt{GPT} & model-name token \\
4 & \texttt{models} & noun or verb depending on context \\
5 & \texttt{learn} & action token \\
6 & \texttt{from} & relation token \\
7 & \texttt{tokens} & object token \\
8 & \texttt{.} & punctuation token \\
\bottomrule
\end{longtable}

The token stream is:

\[
[2,3,4,5,6,7,8].
\]

With \(B=2\), \(T=4\), \(C=6\), and \(V=9\), a batch might be:

\[
x =
\begin{bmatrix}
2 & 3 & 4 & 5\\
3 & 4 & 5 & 6
\end{bmatrix}
\in \mathbb{Z}^{2\times4}.
\]

The embedding layer turns this into:

\[
X \in \mathbb{R}^{2\times4\times6}.
\]

\begin{center}
\begin{tikzpicture}[
  cell/.style={draw, minimum width=0.62cm, minimum height=0.46cm, align=center},
  box/.style={draw, rounded corners=2pt, minimum height=0.65cm, text width=2.2cm, align=center},
  arrow/.style={-{Latex[length=2mm]}, thick},
  note/.style={font=\small, align=center}
]
\node[box, fill=blue!5] (ids) at (0,0) {token IDs\\\shape{[2,4]}};
\foreach \r/\vals in {0/{2,3,4,5},1/{3,4,5,6}} {
  \foreach \c in {0,1,2,3} {
    \pgfmathtruncatemacro{\idx}{\c+1}
    \node[cell, fill=blue!5] at (2.2+\c*0.62,-\r*0.46) {\scriptsize \pgfmathparse{{\vals}[\c]}\pgfmathresult};
  }
}
\node[box, fill=red!8] (lookup) at (5.3,-0.23) {embedding\\lookup table};
\node[box, fill=green!10] (vecs) at (8.4,-0.23) {vectors\\\shape{[2,4,6]}};
\draw[arrow] (ids) -- (1.65,-0.23);
\draw[arrow] (4.95,-0.23) -- (lookup);
\draw[arrow] (lookup) -- (vecs);
\node[note] at (5.3,-1.25) {IDs choose rows; no arithmetic meaning is assumed.};
\end{tikzpicture}
\end{center}

\section{Token IDs Are Names, Not Quantities}

A beginner mistake is to treat token IDs like ordinary numbers. If \texttt{Tiny} has ID \(2\), \texttt{GPT} has ID \(3\), and \texttt{tokens} has ID \(7\), that does not mean:

\[
\texttt{Tiny} < \texttt{GPT} < \texttt{tokens}
\]

in any semantic sense. The IDs are assigned by a tokenizer. They are more like row numbers in a spreadsheet than measurements.

\begin{center}
\begin{tikzpicture}[
  box/.style={draw, rounded corners=2pt, minimum height=0.64cm, text width=2.35cm, align=center},
  arrow/.style={-{Latex[length=2mm]}, thick},
  bad/.style={draw, rounded corners=2pt, minimum height=0.64cm, text width=3.0cm, align=center}
]
\node[box, fill=blue!5] (id2) at (0,0) {ID 2\\\texttt{Tiny}};
\node[box, fill=blue!5] (id3) at (3.0,0) {ID 3\\\texttt{GPT}};
\node[box, fill=blue!5] (id7) at (6.0,0) {ID 7\\\texttt{tokens}};
\node[bad, fill=red!8] (wrong) at (3.0,-1.35) {wrong idea:\\ID distance = meaning distance};
\draw[arrow] (id2) -- (wrong);
\draw[arrow] (id3) -- (wrong);
\draw[arrow] (id7) -- (wrong);
\end{tikzpicture}
\end{center}

The model needs a trainable numerical representation. That is why the first layer of a GPT model is an embedding table.

\section{One-Hot Vectors: The Conceptual Bridge}

Before learning embeddings, it helps to imagine a token as a one-hot vector. For \(V=9\), the token \texttt{Tiny} with ID \(2\) could be represented as:

\[
\mathbf{e}_2 =
\begin{bmatrix}
0&0&1&0&0&0&0&0&0
\end{bmatrix}.
\]

This vector has length \(V\). It says ``select row 2'' and nothing more.

\begin{center}
\begin{tikzpicture}[
  cell/.style={draw, minimum width=0.56cm, minimum height=0.48cm, align=center},
  arrow/.style={-{Latex[length=2mm]}, thick},
  note/.style={font=\small, align=center}
]
\node[note] at (0,0.75) {one-hot vector for ID 2};
\foreach \i in {0,1,2,3,4,5,6,7,8} {
  \ifnum\i=2
    \node[cell, fill=red!12] (c\i) at (\i*0.58,0) {\scriptsize 1};
  \else
    \node[cell, fill=blue!5] (c\i) at (\i*0.58,0) {\scriptsize 0};
  \fi
  \node[font=\tiny] at (\i*0.58,-0.45) {\i};
}
\node[note] (select) at (6.7,0) {selects\\row 2};
\draw[arrow] (c2) -- (select);
\end{tikzpicture}
\end{center}

If \(E_{\text{tok}}\in\mathbb{R}^{V\times C}\), then multiplying a one-hot row vector by the embedding table selects the corresponding row:

\[
\mathbf{e}_i E_{\text{tok}} = E_{\text{tok}}[i].
\]

In practice, PyTorch does not build a huge one-hot vector. It performs a direct lookup, which is faster and uses less memory.

\begin{lstlisting}[language=Python]
tok_emb = nn.Embedding(vocab_size, n_embd)
vec = tok_emb(torch.tensor([2]))   # selects row 2
\end{lstlisting}

\section{The Token Embedding Table}

Let vocabulary size be \(V\) and embedding dimension be \(C\). The token embedding table is:

\[
E_{\text{tok}} \in \mathbb{R}^{V \times C}.
\]

For TinyGPT:

\[
E_{\text{tok}} \in \mathbb{R}^{9 \times 6}.
\]

Each row is a trainable vector. For example:

\[
E_{\text{tok}}[2] =
\begin{bmatrix}
0.12&-0.04&0.31&0.08&-0.17&0.22
\end{bmatrix}.
\]

\begin{center}
\begin{tikzpicture}[
  cell/.style={draw, minimum width=0.62cm, minimum height=0.42cm, align=center},
  rowlab/.style={font=\scriptsize, align=right},
  arrow/.style={-{Latex[length=2mm]}, thick},
  note/.style={font=\small, align=center}
]
\node[note] at (2.5,0.6) {embedding table \(E_{\text{tok}}\in\mathbb{R}^{9\times6}\)};
\foreach \r/\tok in {0/<pad>,1/<unk>,2/Tiny,3/GPT,4/models,5/learn,6/from,7/tokens,8/.} {
  \node[rowlab, anchor=east] at (-0.15,-\r*0.42) {\r\ \texttt{\tok}};
  \foreach \c in {0,1,2,3,4,5} {
    \ifnum\r=2
      \node[cell, fill=red!10] (r\r c\c) at (\c*0.62,-\r*0.42) {};
    \else
      \node[cell, fill=blue!4] at (\c*0.62,-\r*0.42) {};
    \fi
  }
}
\node[note] (id) at (5.1,-0.84) {ID 2 selects\\the \texttt{Tiny} row};
\draw[arrow] (id.west) -- (r2c5.east);
\node[note] at (1.55,-4.15) {The selected row is a length-\(C\) vector used by the model.};
\end{tikzpicture}
\end{center}

\section{Lookup For A Batch}

For a single token ID \(i\), lookup returns one vector:

\[
E_{\text{tok}}[i]\in\mathbb{R}^{C}.
\]

For a sequence of \(T\) IDs, lookup returns a matrix:

\[
E_{\text{tok}}[x_{0:T-1}]\in\mathbb{R}^{T\times C}.
\]

For a batch of \(B\) sequences, lookup returns a rank-3 tensor:

\[
E_{\text{tok}}[x]\in\mathbb{R}^{B\times T\times C}.
\]

\begin{center}
\begin{tikzpicture}[
  box/.style={draw, rounded corners=2pt, minimum height=0.66cm, text width=2.15cm, align=center},
  arrow/.style={-{Latex[length=2mm]}, thick},
  note/.style={font=\small, align=center}
]
\node[box, fill=blue!5] (one) at (0,0) {one ID\\\shape{[]}};
\node[box, fill=red!8] (onev) at (3,0) {one vector\\\shape{[C]}};
\node[box, fill=blue!5] (seq) at (0,-1.15) {one sequence\\\shape{[T]}};
\node[box, fill=red!8] (seqv) at (3,-1.15) {one matrix\\\shape{[T,C]}};
\node[box, fill=blue!5] (batch) at (0,-2.30) {one batch\\\shape{[B,T]}};
\node[box, fill=red!8] (batchv) at (3,-2.30) {one tensor\\\shape{[B,T,C]}};
\draw[arrow] (one) -- (onev);
\draw[arrow] (seq) -- (seqv);
\draw[arrow] (batch) -- (batchv);
\node[note] at (1.5,-3.15) {The same embedding table handles all three cases.};
\end{tikzpicture}
\end{center}

\section{Position Embeddings: Giving Tokens An Address}

Token embeddings identify \term{what} token is present. A GPT model also needs to know \term{where} the token appears.

The sentences:

\begin{quote}
\texttt{models learn}
\end{quote}

and:

\begin{quote}
\texttt{learn models}
\end{quote}

contain the same token IDs but in different positions. Without position information, the initial token vectors would not know their order.

TinyGPT uses a learned position embedding table:

\[
E_{\text{pos}}\in\mathbb{R}^{T_{\max}\times C}.
\]

For a block size of \(T_{\max}=4\) and \(C=6\):

\[
E_{\text{pos}}\in\mathbb{R}^{4\times6}.
\]

At position \(t\), the model adds:

\[
X_{b,t,:}=E_{\text{tok}}[x_{b,t}]+E_{\text{pos}}[t].
\]

\begin{center}
\begin{tikzpicture}[
  cell/.style={draw, minimum width=0.58cm, minimum height=0.46cm, align=center},
  arrow/.style={-{Latex[length=2mm]}, thick},
  plus/.style={circle, draw, minimum size=0.5cm, align=center},
  note/.style={font=\small, align=center}
]
\foreach \i/\v in {0/0.12,1/-0.04,2/0.31,3/0.08,4/-0.17,5/0.22} {
  \node[cell, fill=blue!5] (t\i) at (\i*0.58,0) {\scriptsize \v};
}
\node[note, anchor=east] at (-0.25,0) {token vector\\\texttt{Tiny}};
\foreach \i/\v in {0/0.01,1/0.05,2/-0.02,3/0.03,4/0.04,5/-0.01} {
  \node[cell, fill=yellow!18] (p\i) at (\i*0.58,-0.8) {\scriptsize \v};
}
\node[note, anchor=east] at (-0.25,-0.8) {position vector\\\(t=0\)};
\node[plus] (add) at (4.1,-0.4) {\(+\)};
\foreach \i/\v in {0/0.13,1/0.01,2/0.29,3/0.11,4/-0.13,5/0.21} {
  \node[cell, fill=red!8] (x\i) at (5.0+\i*0.58,-0.4) {\scriptsize \v};
}
\node[note] at (6.45,0.45) {initial representation\\for \texttt{Tiny} at position 0};
\draw[arrow] (t5) -- (add);
\draw[arrow] (p5) -- (add);
\draw[arrow] (add) -- (x0);
\end{tikzpicture}
\end{center}

The addition works because both vectors have the same length \(C\). The result still has shape \shape{[C]}; it simply carries two kinds of information.

\section{Broadcasting Position Embeddings Across The Batch}

In source code, token embeddings have shape \shape{[B,T,C]}. Position embeddings for the positions \(0,1,\ldots,T-1\) have shape \shape{[T,C]}. PyTorch broadcasts the position embeddings across the batch axis:

\[
\shape{[B,T,C]} + \shape{[T,C]} \rightarrow \shape{[B,T,C]}.
\]

\begin{center}
\begin{tikzpicture}[
  box/.style={draw, rounded corners=2pt, minimum height=0.65cm, text width=2.45cm, align=center},
  arrow/.style={-{Latex[length=2mm]}, thick},
  plus/.style={circle, draw, minimum size=0.5cm, align=center},
  note/.style={font=\small, align=center}
]
\node[box, fill=blue!5] (tok) at (0,0) {token embeddings\\\shape{[B,T,C]}};
\node[box, fill=yellow!18] (pos) at (0,-1.0) {position embeddings\\\shape{[T,C]}};
\node[plus] (add) at (3.1,-0.5) {\(+\)};
\node[box, fill=red!8] (x) at (6.1,-0.5) {initial hidden state\\\shape{[B,T,C]}};
\draw[arrow] (tok) -- (add);
\draw[arrow] (pos) -- (add);
\draw[arrow] (add) -- (x);
\node[note] at (3.1,-1.55) {same position table is reused for every batch example};
\end{tikzpicture}
\end{center}

In PyTorch:

\begin{lstlisting}[language=Python]
B, T = idx.shape
pos = torch.arange(0, T, device=idx.device)  # [T]
tok = self.tok_emb(idx)                      # [B, T, C]
pos = self.pos_emb(pos)                      # [T, C]
x = tok + pos                                # [B, T, C]
\end{lstlisting}

The line \texttt{x = tok + pos} is small, but conceptually it creates the first representation the Transformer will process.

\section{Representations Are Working States}

An embedding vector is not a dictionary definition. It is a working state used by the model. The vector for \texttt{models} at the input may mostly encode token identity and position. After attention layers, the vector at that position may include information from earlier tokens.

\begin{center}
\begin{tikzpicture}[
  box/.style={draw, rounded corners=2pt, minimum height=0.65cm, text width=2.25cm, align=center},
  arrow/.style={-{Latex[length=2mm]}, thick},
  note/.style={font=\small, align=center}
]
\node[box, fill=blue!5] (e0) at (0,0) {input embedding\\identity + position};
\node[box, fill=red!8] (e1) at (3,0) {after block 1\\local context};
\node[box, fill=red!8] (e2) at (6,0) {after block 2\\richer context};
\node[box, fill=green!10] (logit) at (9,0) {output head\\vocabulary scores};
\draw[arrow] (e0) -- (e1);
\draw[arrow] (e1) -- (e2);
\draw[arrow] (e2) -- (logit);
\node[note] at (4.5,-0.9) {The vector keeps the same width \(C\), but its information changes.};
\end{tikzpicture}
\end{center}

This is why the word \term{representation} is useful. A representation is not only a vector. It is a vector with a role at a particular stage of computation.

\section{How Embeddings Learn}

At initialization, embedding rows are random. During training, the loss creates gradients that update only the rows used by the batch.

Suppose the context is:

\[
\texttt{Tiny GPT models learn from}
\]

and the target next token is:

\[
\texttt{tokens}.
\]

If the model assigns too little probability to \texttt{tokens}, the cross-entropy loss is high. Backpropagation sends gradients through:

\[
\text{loss}\rightarrow \text{logits}\rightarrow \text{Transformer blocks}\rightarrow \text{input embeddings}.
\]

\begin{center}
\begin{tikzpicture}[
  box/.style={draw, rounded corners=2pt, minimum height=0.65cm, text width=2.25cm, align=center},
  arrow/.style={-{Latex[length=2mm]}, thick},
  back/.style={{Latex[length=2mm]}-, thick, red!70!black},
  note/.style={font=\small, align=center}
]
\node[box, fill=blue!5] (rows) at (0,0) {selected rows\\IDs 2,3,4,5};
\node[box, fill=red!8] (model) at (3,0) {TinyGPT\\forward pass};
\node[box, fill=green!10] (logits) at (6,0) {logits for\\next tokens};
\node[box, fill=yellow!18] (loss) at (9,0) {loss};
\draw[arrow] (rows) -- (model);
\draw[arrow] (model) -- (logits);
\draw[arrow] (logits) -- (loss);
\draw[back] (rows.south) .. controls (3,-1.2) and (6,-1.2) .. (loss.south);
\node[note] at (4.5,-1.55) {Gradients adjust embedding rows that participated in the example.};
\end{tikzpicture}
\end{center}

In source code, the embedding table is an \texttt{nn.Parameter} inside \texttt{nn.Embedding}. The optimizer updates it like any other trainable weight.

\section{Geometry: What Similarity Can Mean}

After training, embeddings used in similar contexts may point in similar directions. A common measure is cosine similarity:

\[
\cos(\mathbf{a},\mathbf{b})
=
\frac{\mathbf{a}\cdot\mathbf{b}}
{\|\mathbf{a}\|\,\|\mathbf{b}\|}.
\]

Cosine similarity is \(1\) when two vectors point in the same direction, \(0\) when they are orthogonal, and \(-1\) when they point in opposite directions.

\begin{center}
\begin{tikzpicture}[
  arrow/.style={-{Latex[length=2mm]}, thick},
  note/.style={font=\small, align=center}
]
\draw[->] (-0.2,0) -- (5.4,0) node[right] {\scriptsize feature 1};
\draw[->] (0,-0.2) -- (0,3.2) node[above] {\scriptsize feature 2};
\draw[arrow, blue!70!black] (0,0) -- (2.2,2.0) node[above] {\scriptsize \texttt{Tiny}};
\draw[arrow, blue!70!black] (0,0) -- (2.7,2.2) node[right] {\scriptsize \texttt{GPT}};
\draw[arrow, red!70!black] (0,0) -- (4.3,0.7) node[right] {\scriptsize \texttt{.}};
\draw[dashed] (2.2,2.0) -- (2.7,2.2);
\node[note] at (2.8,-0.65) {This 2D picture is only an illustration. Real embeddings may have hundreds or thousands of dimensions.};
\end{tikzpicture}
\end{center}

Be careful: embedding geometry is useful, but it is not magic. The model learns whatever geometry helps reduce next-token loss on the training data.

\section{The Output Head And Weight Tying}

The embedding table reads tokens into the model:

\[
E_{\text{tok}}\in\mathbb{R}^{V\times C}.
\]

The language-model head maps hidden vectors back to vocabulary logits:

\[
z = hW_{\text{out}} + b,
\qquad
W_{\text{out}}\in\mathbb{R}^{C\times V}.
\]

Many GPT models tie the input embedding table and output head:

\[
W_{\text{out}} = E_{\text{tok}}^\top.
\]

Conceptually:

\begin{itemize}
  \item input side: ``What vector should this token ID become?''
  \item output side: ``Which token vector is most compatible with this hidden state?''
\end{itemize}

\begin{center}
\begin{tikzpicture}[
  box/.style={draw, rounded corners=2pt, minimum height=0.65cm, text width=2.45cm, align=center},
  arrow/.style={-{Latex[length=2mm]}, thick},
  note/.style={font=\small, align=center}
]
\node[box, fill=blue!5] (ids) at (0,0) {input IDs};
\node[box, fill=red!8] (emb) at (3,0) {token embedding\\\(E_{\text{tok}}\)};
\node[box, fill=green!10] (hidden) at (6,0) {hidden state\\\(h\)};
\node[box, fill=red!8] (head) at (9,0) {output head\\\(E_{\text{tok}}^\top\)};
\node[box, fill=yellow!18] (logits) at (12,0) {logits};
\draw[arrow] (ids) -- (emb);
\draw[arrow] (emb) -- (hidden);
\draw[arrow] (hidden) -- (head);
\draw[arrow] (head) -- (logits);
\draw[<->, thick, dashed] (emb.south) .. controls (5,-1.3) and (7,-1.3) .. (head.south);
\node[note] at (6,-1.65) {weight tying reuses the same token table for reading and scoring};
\end{tikzpicture}
\end{center}

In PyTorch, weight tying can be written as:

\begin{lstlisting}[language=Python]
self.tok_emb = nn.Embedding(vocab_size, n_embd)
self.lm_head = nn.Linear(n_embd, vocab_size, bias=False)
self.lm_head.weight = self.tok_emb.weight
\end{lstlisting}

The shapes work because \texttt{nn.Linear(C,V)} stores its weight as \shape{[V,C]}. That is the same shape as \texttt{tok\_emb.weight}.

\section{Parameter Count}

The token embedding table has:

\[
V\times C
\]

parameters. The position embedding table has:

\[
T_{\max}\times C
\]

parameters.

For TinyGPT:

\[
V=9,\qquad T_{\max}=4,\qquad C=6.
\]

So:

\[
\text{token embedding params}=9\times6=54,
\]

\[
\text{position embedding params}=4\times6=24.
\]

Total embedding parameters:

\[
54+24=78.
\]

\begin{center}
\begin{tikzpicture}[
  box/.style={draw, rounded corners=2pt, minimum height=0.68cm, text width=2.7cm, align=center},
  arrow/.style={-{Latex[length=2mm]}, thick}
]
\node[box, fill=blue!5] (tok) at (0,0) {token embeddings\\\(V C = 54\)};
\node[box, fill=yellow!18] (pos) at (3.4,0) {position embeddings\\\(T_{\max}C = 24\)};
\node[box, fill=red!8] (total) at (6.8,0) {embedding total\\78 parameters};
\draw[arrow] (tok) -- (pos);
\draw[arrow] (pos) -- (total);
\end{tikzpicture}
\end{center}

For a larger GPT-style model with \(V=50{,}257\) and \(C=768\):

\[
50{,}257\times768=38{,}597{,}376
\]

token embedding parameters. This is why vocabulary size matters. A larger vocabulary may shorten token sequences, but it also expands the embedding table and output vocabulary.

\section{Embedding Source Code Logic}

A minimal TinyGPT embedding module looks like this:

\begin{lstlisting}[language=Python]
class TinyGPTEmbeddings(nn.Module):
    def __init__(self, vocab_size, block_size, n_embd):
        super().__init__()
        self.vocab_size = vocab_size
        self.block_size = block_size
        self.n_embd = n_embd
        self.tok_emb = nn.Embedding(vocab_size, n_embd)
        self.pos_emb = nn.Embedding(block_size, n_embd)

    def forward(self, idx):
        # idx: [B, T]
        B, T = idx.shape
        assert T <= self.block_size
        assert idx.min() >= 0
        assert idx.max() < self.vocab_size

        pos = torch.arange(0, T, device=idx.device)  # [T]
        tok = self.tok_emb(idx)                      # [B, T, C]
        pos = self.pos_emb(pos)                      # [T, C]
        x = tok + pos                                # [B, T, C]
        return x
\end{lstlisting}

The assertions are educational but valuable. They catch the most common data/model mismatch bugs:

\begin{itemize}
  \item token ID outside the vocabulary;
  \item sequence length longer than block size;
  \item tokenizer vocabulary not matching model configuration.
\end{itemize}

\begin{center}
\begin{tikzpicture}[
  box/.style={draw, rounded corners=2pt, minimum height=0.62cm, text width=2.3cm, align=center},
  arrow/.style={-{Latex[length=2mm]}, thick},
  note/.style={font=\small, align=center}
]
\node[box, fill=blue!5] (idx) at (0,0) {\texttt{idx}\\\shape{[B,T]}};
\node[box, fill=gray!12] (check) at (2.8,0) {shape and ID\\checks};
\node[box, fill=blue!5] (tok) at (5.6,0.65) {\texttt{tok\_emb}\\\shape{[B,T,C]}};
\node[box, fill=yellow!18] (pos) at (5.6,-0.65) {\texttt{pos\_emb}\\\shape{[T,C]}};
\node[box, fill=red!8] (x) at (8.7,0) {initial state\\\shape{[B,T,C]}};
\draw[arrow] (idx) -- (check);
\draw[arrow] (check) -- (tok);
\draw[arrow] (check) -- (pos);
\draw[arrow] (tok) -- (x);
\draw[arrow] (pos) -- (x);
\node[note] at (4.8,-1.55) {This module prepares the input for the first Transformer block.};
\end{tikzpicture}
\end{center}

\section{Engineering Practice}

Embedding code should check:

\begin{itemize}
  \item token IDs are within \([0,V-1]\);
  \item sequence length does not exceed the configured block size;
  \item embedding output shape is \shape{[B,T,C]};
  \item tokenizer vocabulary size matches model config;
  \item checkpoint metadata records tokenizer identity;
  \item tied or untied output-head choice is saved in config;
  \item device placement is consistent for input IDs and embedding parameters.
\end{itemize}

For the \texttt{llm-train} project, the tokenizer and model config should be treated as one contract. If a checkpoint was trained with one vocabulary, loading it with a different tokenizer is usually invalid even if the code runs.

\section{Hands-On Lab: Inspect TinyGPT Embeddings}

Create a small script that builds an embedding module and traces shapes:

\begin{lstlisting}[language=Python]
import torch
import torch.nn as nn

torch.manual_seed(1)
B, T, C, V = 2, 4, 6, 9

idx = torch.tensor([[2, 3, 4, 5],
                    [3, 4, 5, 6]], dtype=torch.long)

tok_emb = nn.Embedding(V, C)
pos_emb = nn.Embedding(T, C)

pos = torch.arange(0, T)
tok = tok_emb(idx)
pos_vec = pos_emb(pos)
x = tok + pos_vec

print("idx:", idx.shape)
print("tok:", tok.shape)
print("pos_vec:", pos_vec.shape)
print("x:", x.shape)
print("Tiny row:", tok_emb.weight[2])
\end{lstlisting}

Then add a language-model head and tie weights:

\begin{lstlisting}[language=Python]
lm_head = nn.Linear(C, V, bias=False)
lm_head.weight = tok_emb.weight

logits = lm_head(x)
print("logits:", logits.shape)
print("shared object:", lm_head.weight is tok_emb.weight)
\end{lstlisting}

Expected shapes:

\begin{lstlisting}
idx:     [2, 4]
tok:     [2, 4, 6]
pos_vec: [4, 6]
x:       [2, 4, 6]
logits:  [2, 4, 9]
\end{lstlisting}

\checkpoint

\begin{enumerate}
  \item Why are token IDs not meaningful numerical distances?
  \item What shape is the TinyGPT token embedding table when \(V=9\) and \(C=6\)?
  \item Why does a sequence of IDs with shape \shape{[T]} become \shape{[T,C]} after embedding?
  \item Why does GPT need position embeddings?
  \item What does \texttt{x = tok + pos} combine?
  \item What is weight tying?
  \item Why must tokenizer identity be stored with a checkpoint?
  \item If \(V=8192\), \(T_{\max}=256\), and \(C=384\), how many token and position embedding parameters are there?
\end{enumerate}

\subsection*{Solutions}

\begin{enumerate}
  \item Token IDs are labels assigned by the tokenizer. Their numeric distance does not encode semantic distance.
  \item The table has shape \shape{[9,6]} because it has one row for each vocabulary token and six features per row.
  \item Each ID selects one length-\(C\) row from the embedding table. A length-\(T\) list of IDs selects \(T\) rows, giving \shape{[T,C]}.
  \item Token embeddings alone do not tell the model where a token appears. Position embeddings give each token representation an address in the context window.
  \item It combines token identity information and position information into the initial hidden state.
  \item Weight tying reuses the token embedding table as the output vocabulary classifier, usually by sharing \texttt{tok\_emb.weight} with \texttt{lm\_head.weight}.
  \item The rows of the embedding table correspond to token IDs. If the tokenizer changes, row meanings change, so the checkpoint can be misinterpreted.
  \item Token embeddings: \(8192\times384=3{,}145{,}728\). Position embeddings: \(256\times384=98{,}304\). Total: \(3{,}244{,}032\).
\end{enumerate}
